\begin{table}[htbp]
\caption{Uttar Pradesh: Electoral Balance Tests for Quota Assignment}
\label{tab:balance_electoral_up}
{\centering\scriptsize
\begin{tabular}{@{}lrrr rrr rrr@{}}
\toprule
& \multicolumn{3}{c}{2005$\rightarrow$2010} & \multicolumn{3}{c}{2010$\rightarrow$2015} & \multicolumn{3}{c}{2015$\rightarrow$2021} \\
\cmidrule(lr){2-4} \cmidrule(lr){5-7} \cmidrule(lr){8-10}
Variable & Quota & Open & p & Quota & Open & p & Quota & Open & p \\
\midrule
Prior female winner & 0.51 & 0.53 & 0.000 & 0.36 & 0.45 & 0.000 & 0.29 & 0.52 & 0.000 \\
OBC reservation & 0.21 & 0.30 & 0.000 & 0.26 & 0.29 & 0.000 & 0.17 & 0.32 & 0.000 \\
SC/ST reservation & 0.22 & 0.21 & 0.000 & 0.15 & 0.25 & 0.000 & 0.27 & 0.18 & 0.000 \\
\midrule
N & \multicolumn{3}{c}{39,510} & \multicolumn{3}{c}{34,556} & \multicolumn{3}{c}{41,924} \\
\bottomrule
\end{tabular}
\par}

\vspace{0.5ex}
\parbox{\linewidth}{\scriptsize \emph{Notes:} Balance tests for quota assignment on prior electoral characteristics in Uttar Pradesh. ``Quota'' = GPs reserved for women in the later year; ``Open'' = GPs not reserved. Prior female winner = proportion of GPs with a female pradhan in the prior election. OBC/SC/ST reservation = proportion of GPs with caste-based reservations. p-values from t-tests. This table uses the full electoral panel (not SHRUG-filtered).}
\end{table}
